\section{A scalar core response for cliques with unequal pendants}
\label{sec:op3-local-clique-pendants}

\paragraph{Status and promised family.}
This \statusproved{} draft, awaiting independent review, gives a different
representation for a family on which explicit global refresh is expensive.
The graph is promised to consist of a clique $C$ of unknown size $k\geq3$
with $t_i\geq0$ private degree-one leaves attached to each core vertex $i$.
The physical seed $v$ belongs to the core. Pendant counts can be unequal
and arbitrarily large. The whole scalar-sweep algorithm is implemented;
it does not import a dynamic tree or numerical SDD solver. This is a
structural result, not a graph-uniform OP3 theorem.
The graph promise is an assumption; verifying it on unread rows is not
part of the claimed work.

Fix $0<\alpha<1$ and $\lambda>0$, with the same $\gamma,\bar\alpha$,
physical response and original residual conventions as before.
If $\lambda d_v\geq1$, return zero. Otherwise read the seed row once and
query its neighbors' original degrees. Its neighbors of degree greater
than one, together with $v$, are exactly $C$. Thus $k$, all original core
degrees and all $t_i=d_i-(k-1)$ are discovered with $O(d_v)$ work, without
reading any other core row or an inactive attachment.

\begin{lemma}[Three-piece scalar core response; \statusproved]
\label{lem:op3-clique-scalar-response}
Set $S=\sum_{i\in C}u_i$, $b_i=\mathbb 1_{\{i=v\}}-\lambda d_i$,
$a_i=d_i+\gamma$, and $c_i=a_i-\gamma^2t_i$. The exact obstacle response
is determined by the unique nonnegative fixed point
\begin{equation}
 S=F(S):=\sum_{i\in C}\phi_i(S),\qquad
 \phi_i(S)=
 \begin{cases}
 0,&b_i+\gamma S\leq0,\\
 (b_i+\gamma S)/a_i,&0<b_i+\gamma S\leq a_i\lambda/\gamma,\\
 (b_i+\gamma S-\gamma\lambda t_i)/c_i,
       &b_i+\gamma S>a_i\lambda/\gamma.
 \end{cases}
 \label{eq:op3-clique-response}
\end{equation}
Each pendant at core vertex $i$ then has value
$\max\{0,\gamma\phi_i(S)-\lambda\}$.
The function $F$ is continuous and piecewise affine with at most $2k$
breakpoints. Every affine slope is bounded above by
\begin{equation}
 \frac{k\gamma}{k-1+\gamma}<1.
 \label{eq:op3-clique-slope-bound}
\end{equation}
\end{lemma}
\begin{proof}
A degree-one pendant's obstacle equation gives exactly the displayed
positive-part response. Substituting all $t_i$ identical pendant responses
into the original core equation gives
\[
 (d_i+\gamma)u_i-\gamma t_i(\gamma u_i-\lambda)_+
     =b_i+\gamma S
\]
at a positive core coordinate, with the corresponding nonpositive gate
when it is zero. Solving its two positive linear pieces gives
\cref{eq:op3-clique-response}. Both transitions are continuous and their
slopes are $0$, $\gamma/a_i$ and $\gamma/c_i$. Since
\[
 c_i=k-1+\gamma+(1-\gamma^2)t_i\geq k-1+\gamma>0,
\]
summing the largest slopes gives \cref{eq:op3-clique-slope-bound}.
Thus $F(S)-S$ is strictly decreasing. If $\lambda d_v\geq1$ its value
at zero is zero; otherwise it is positive there and tends to a negative
value as $S$ increases without bound. This proves a unique nonnegative
fixed point. The resulting point satisfies every original core and pendant
obstacle equation or inequality, so strict convexity identifies it as the
unique obstacle solution.
\end{proof}

\paragraph{Exact breakpoint sweep.}
The possible breakpoints of row $i$ are
\[
 S_{i,0}=-b_i/\gamma,\qquad
 S_{i,1}=(a_i\lambda/\gamma-b_i)/\gamma.
\]
The second has no effect when $t_i=0$. Determine every row's affine
piece at $S=0$, ignore its negative or irrelevant breakpoints, and sort
the remaining at most $2k$ events. Keep only the aggregate coefficients
$F(S)=AS+B$ for the current interval. Its candidate root is $B/(1-A)$.
If this is at or before the next breakpoint, it is the fixed point.
Otherwise advance through that event and change only the affected row's
two aggregate coefficients. Continuity and the slope bound guarantee
that the root stays beyond the current interval's left endpoint until it
is found. Ties need no sign margin: a root at a breakpoint is valid in
both neighboring pieces. This is an exact sweep, not a contraction
iteration with a hidden $1/(1-A)$ iteration count.

\begin{theorem}[A local clique-pendant solver; \statusproved]
\label{thm:op3-local-clique-pendants}
On the promised family, the preceding algorithm returns the exact
$\lambda$-obstacle response with
\begin{equation}
 O\!\left(1+\operatorname{cvol}(U)
                  \log(2+\operatorname{cvol}(U))\right)
 \label{eq:op3-clique-local-work}
\end{equation}
local exact-real word work and $O(1+\operatorname{cvol}(U))$ space, where
$U$ is its positive support. All original graph accesses and output are
included. For $0<\epsappr,\rho<1$, $\lambda=\epsappr/2$ gives
$\widetilde O(1/\epsappr)$ ACL work, and, separately, $\lambda=\rho$
gives exact $\widetilde O(1/\rho)$ RPPR work on this family.
\end{theorem}
\begin{proof}
The zero case uses one degree query. In the nonzero case the seed is
positive, and its one scanned row discovers all $k\leq d_v+1$ core
identifiers and their degrees. The scalar sweep stores and sorts at most
$2k$ constant-size records, evaluates at most $2k+1$ candidate roots, and
changes constant-size aggregates at each crossed event. Its work is
$O(k\log(2+k))$ and its additional state is $O(k)$, regardless of the
magnitude of any inactive pendant count.

Evaluate all core values once. Read an additional core row only when its
computed exact value is positive, using the cached seed row for $v$.
For each such row, the core-membership dictionary distinguishes the
private pendants. Query their degrees and output a pendant only when
$\gamma u_i-\lambda>0$; scan that positive degree-one row once if an
explicit original incidence check is desired. Never read a zero core or
zero pendant row. These operations cost $O(\operatorname{vol}(U))$,
including cached incidence reads and every positive output coordinate.
All core-value writes, even at zero core coordinates, cost only $O(k)$
and are charged to the already scanned positive seed row.

The preceding lemma gives the exact original obstacle certificate. Its
mass identity yields $\lambda\operatorname{vol}(U)<1$ when nonzero.
Original full degrees, all discoveries, dictionaries and temporary rows
fit in $O(1+\operatorname{cvol}(U))$ space. Combining these bounds proves
\cref{eq:op3-clique-local-work} and the two distinct accuracy-namespace
consequences. There is no polynomial arithmetic-operation dependence on
$1/\alpha$; rational bit complexity and numerical stability are not claimed.
\end{proof}

\paragraph{Why this complements the refresh obstruction.}
The family of \cref{prop:op3-clique-leaf-refresh} lies in this promised
class. The scalar sweep handles its $\Theta(k^2)$ cycle rank without a
rank factor, whereas the explicit refresh policy must write $\Omega(k^3)$
port coordinates. The useful parameter there is the simple algebraic form
of the core coupling, not its number of independent cycles. This does not
extend automatically to arbitrary cores: several coupled scalar aggregates
would create a multidimensional response and a new event-location problem.

\paragraph{Measured.}
\texttt{local\_clique\_pendants.py} compares its actual local scalar sweep
against independent exact original-matrix obstacle solves for unequal
pendant counts and every core seed. Exact core-birth, pendant-birth and
zero-seed equalities are included. Implicit finite graphs with enormous
inactive pendant counts are checked using original core equations and
pendant inequalities with explicit multiplicities; no dense ambient graph
is built. The audit is saved in \texttt{LOCAL\_CLIQUE\_PENDANT\_AUDIT.json}.
This implementation is separate from the conditional source interfaces
in the general cycle-rank construction.
